% ============================================================================
%  CO3133 Assignment 3, final presentation (M3). Group G-roundhog.
%  Build:  ./checkslides.sh slides.tex      (lualatex, needs slidekit.sty)
%  Speaker copy with notes: \setbeameroption{show notes} in slidekit.sty.
% ============================================================================
\documentclass[8pt,aspectratio=169]{beamer}
\usepackage{slidekit}
\usepackage{tikz}
\usetikzlibrary{positioning,arrows.meta}

\title{Multimodal Recipe Classification}
\subtitle{UPMC Food-101: one web image and its page title, 101 dishes}
\author{Phan Bá Thanh \quad Trần Công Hoàng Phước \quad\gray{Group G-roundhog}}
\date{\gray{CO3133 Deep Learning and Its Applications, Semester 261, Assignment 3 (M3)}}

\begin{document}

\titleframe

\begin{frame}{Outline}
  \renewcommand{\arraystretch}{2.4}
  \makebox[\textwidth]{%
  \begin{tabular}{@{}r@{\hspace{1.4em}}l@{\hspace{2.4em}}l@{}}
    \gray{1} & Problem and data & $79\%$ of titles name the class \\
    \gray{2} & Models           & UA, UB, F0, F1 and the new \K{F2} \\
    \gray{3} & Results          & F2 $>$ F1, F2 $\approx$ F0, F2 best on masked titles \\
    \gray{4} & Analysis         & conflicts, attention, error taxonomy \\
  \end{tabular}}
\end{frame}

% ===========================================================================
\section{Problem and data}
% ===========================================================================

\begin{frame}{Task}
  \takeaway{Input: one web image + the page title. Output: one of 101 dishes.}

  \fig{pairs}

  \showtakeaway
  \note{UPMC Food-101, Kaggle mirror with titles only. Complete pairs: 62,771 train,
    4,824 val, 22,594 test. Image and title come from the same recipe page. Split grouped by
    image-duplicate clusters. Manual audit of 200 pairs: 81.5 \% valid, 0.5 \% wrong labels.}
\end{frame}

\begin{frame}{Two leaks}
  \takeaway{The title often names the label: every claim is also checked on masked and unseen titles.}

  \makebox[\textwidth]{%
  \slidetable
  \begin{tabular}{@{}llr@{}}
    \toprule
    \headrow
    Leak & Effect & Size \\
    \midrule
    class name in the title & ``\textless class\textgreater{} recipe'' search query & $79\%$ of titles \\
    same title in train and test & several images of one page & $39\%$ of test titles \\
    same image in train and test & near-duplicate photos & $14.5\%$ of test images \\
    \midrule
    \K{Control} & E2 masks class names; \emph{title-unseen} and \emph{dedup} test views & \\
    \bottomrule
  \end{tabular}}

  \showtakeaway
  \note{Keyword rule K0 alone: 77 \% accuracy. Title-unseen view: 13,701 test items whose title
    never occurs in train or val. Dedup view: 19,270 items without a train-linked image.}
\end{frame}

% ===========================================================================
\section{Models}
% ===========================================================================

\begin{frame}{Models}
  \takeaway{Two unimodal baselines, three fusions, two text references; 3 seeds each.}

  \makebox[\textwidth]{%
  \slidetable
  \begin{tabular}{@{}lll@{}}
    \toprule
    \headrow
    ID & Model & Fusion \\
    \midrule
    UA & ResNet-50, image only & none \\
    UB & DistilBERT, title only & none \\
    K0 / T0 & keyword rule / TF-IDF + logistic regression & none \\
    F0 & late: $w\,p_{\text{UA}} + (1-w)\,p_{\text{UB}}$ & output averaging \\
    F1 & concat of two pooled vectors $\to$ MLP & intermediate \\
    \K{F2} & cross-modal transformer + aux.\ losses + modality dropout & token level \\
    \bottomrule
  \end{tabular}}

  \showtakeaway
  \note{All models: same split, same protocol. Seeds 42, 123, 2026. Selection: best val
    macro-F1. Decision rule: mean macro-F1 difference larger than the larger seed std AND
    McNemar p < 0.05 on seed 42.}
\end{frame}

\begin{frame}{Why a new fusion}
  \takeaway{M2: F1 lets the title do all the work, and its image branch stays weak.}

  \makebox[\textwidth]{%
  \slidetable
  \begin{tabular}{@{}lrr@{}}
    \toprule
    \headrow
    Accuracy (\%) & F1 & reference \\
    \midrule
    training set & $99.7$ & \\
    title removed at test time & \F{$53.3$} & UA image only: $70.7$ \\
    only the image model right (Conflict-I) & \F{$72.1$} & F0: $92.7$ \\
    \bottomrule
  \end{tabular}}

  \showtakeaway
  \note{Modality laziness (Wang et al. 2020): the title alone fits the training labels, so the
    joint loss stops pushing the image encoder. F2 is built to fix exactly this.}
\end{frame}

\begin{frame}{F2: cross-modal fusion}
  \takeaway{Region and word tokens attend to each other; aux.\ losses and modality dropout keep both encoders working.}

  \dia{%
  \begin{tikzpicture}[every node/.style={font=\normalsize},
      box/.style={draw=key, rounded corners, align=center, minimum height=9mm, inner sep=4pt, fill=white},
      aux/.style={draw=foil, dashed, rounded corners, align=center, inner sep=3pt},
      >={Stealth}]
    \node[box] (img) at (0,0) {image};
    \node[box] (rn)  at (2.8,0) {ResNet-50\\49 regions};
    \node[box] (txt) at (0,-2.4) {title};
    \node[box] (db)  at (2.8,-2.4) {DistilBERT\\word tokens};
    \node[box, fill=paper, minimum height=26mm, minimum width=30mm] (tf) at (7.2,-1.2) {2 transformer\\layers};
    \node[box] (fus) at (7.2,1.0) {\texttt{[FUS]} token};
    \node[box] (out) at (10.6,-1.2) {101 dishes};
    \node[aux] (ai) at (2.8,1.4) {aux.\ image head};
    \node[aux] (at) at (2.8,-3.8) {aux.\ title head};
    \node[aux] (md) at (7.2,-3.6) {15\%: drop one modality\\$\to$ \texttt{[MISSING]}};
    \draw[->] (img) -- (rn); \draw[->] (txt) -- (db);
    \draw[->] (rn) -- (tf.west |- rn); \draw[->] (db) -- (tf.west |- db);
    \draw[->] (fus) -- (tf); \draw[->] (tf) -- (out);
    \draw[->, foil, dashed] (rn) -- (ai); \draw[->, foil, dashed] (db) -- (at);
    \draw[->, foil, dashed] (md) -- (tf);
  \end{tikzpicture}}

  \showtakeaway
  \note{Image tokens: ResNet-50 layer4, 7x7 = 49 regions, projected to 512, learned position.
    Title tokens: DistilBERT hidden states projected to 512. Two pre-LN layers, 8 heads.
    Loss = CE(fused) + 0.5 (CE(image head) + CE(title head)). Epoch 1 with frozen encoders.
    98.0 M parameters, 5.06 GMACs per item (F1: 4.73).}
\end{frame}

\begin{frame}{Controlled experiments}
  \takeaway{Each experiment changes one factor and keeps the split and the seeds.}

  \makebox[\textwidth]{%
  \slidetable
  \begin{tabular}{@{}lll@{}}
    \toprule
    \headrow
    ID & Change & Question \\
    \midrule
    E1 & no auxiliary losses ($\lambda = 0$) & do they keep the image branch alive? \\
    E4 & no modality dropout ($p = 0$) & robustness to a missing modality? \\
    E2 & class names $\to$ \texttt{[MASK]} in all titles & what is left without the label leak? \\
    E-mod & remove title or image from $0$--$100\%$ of test pairs & how fast does each model degrade? \\
    \bottomrule
  \end{tabular}}

  \showtakeaway
  \note{E1: aux heads still trained on detached features, as linear probes, so we can measure
    the encoders. E2: retrain UB, T0, F0, F1, F2 on masked titles; 79.8 \% of train titles change,
    0 residual mentions. 18 new runs on the lab RTX A6000, about 5.7 GPU-hours.}
\end{frame}

% ===========================================================================
\section{Results}
% ===========================================================================

\begin{frame}{Main result}
  \takeaway{Every fusion beats every single modality by 6 to 7 points of macro-F1.}

  \fig{macro_f1}

  \showtakeaway
  \note{Test complete pairs, 22,594 items, mean and std over 3 seeds. Same ranking on the dedup
    and title-unseen views. T0 slightly beats UB: short titles that mostly name the dish do not
    need a transformer.}
\end{frame}

\begin{frame}{Decision rule}
  \takeaway{F2 beats F1, but ties the late fusion F0.}

  \makebox[\textwidth]{%
  \slidetable
  \begin{tabular}{@{}lrrrl@{}}
    \toprule
    \headrow
    A vs B & $\Delta$ macro-F1 & larger std & McNemar $p$ & verdict \\
    \midrule
    F2 vs UB & $+6.73$ & $0.25$ & $< 10^{-300}$ & F2 better \\
    F1 vs F0 & $-1.06$ & $0.24$ & $1.5\cdot10^{-8}$ & F0 better \\
    F2 vs F1 & $+0.73$ & $0.07$ & $2.3\cdot10^{-9}$ & \K{F2 better} \\
    F2 vs F0 & $-0.33$ & $0.24$ & $0.50$ & \F{no decision} \\
    \bottomrule
  \end{tabular}}

  \showtakeaway
  \note{Rule: difference of mean macro-F1 larger than the larger seed std, and McNemar p < 0.05 on
    seed 42. F2 vs F0: 465 vs 487 discordant items. With two strong unimodal models and
    informative titles, averaging probabilities is a hard baseline. F2 is also the best
    calibrated model: ECE 2.9 \%, 0.45 \% after temperature scaling; F0 is 16.0 \%.}
\end{frame}

\begin{frame}{E1: auxiliary losses}
  \takeaway{Without the auxiliary losses the image encoder weakens: $-7.4$ points on its own.}

  \makebox[\textwidth]{%
  \slidetable
  \begin{tabular}{@{}lrrr@{}}
    \toprule
    \headrow
    & macro-F1 & image head acc. & acc.\ title removed \\
    \midrule
    F2 ($\lambda = 0.5$) & $93.53$ & \K{$69.7$} & $65.3$ \\
    E1 ($\lambda = 0$)   & $93.07$ & $62.3$ & $58.6$ \\
    \midrule
    \gray{UA image only} & \gray{$70.8$} & \gray{$70.7$} & \\
    \bottomrule
  \end{tabular}}

  \showtakeaway
  \note{Decision rule: F2 better, +0.46 pp, std 0.08, McNemar p = 3e-7. The image head of F2 is as
    good as the standalone ResNet-50; the linear probe on E1's image features is 7 points worse.
    Text head: 87.1 vs 76.8. This is the modality laziness of F1, measured directly.}
\end{frame}

\begin{frame}{Missing title}
  \takeaway{Aux.\ losses and modality dropout each add 5 to 7 points; F0 still degrades best.}

  \fig{removal_title}

  \showtakeaway
  \note{Title removed from a random fraction of test pairs, no retraining. Lines are linear by
    construction; read the end points. 100 \% removed: F0 70.8 (falls back to UA), F2 65.3,
    E4 60.1, E1 58.5, F1 52.5. E4 vs F2 on complete pairs: +0.08, no decision, so modality dropout
    is free. Removing the image instead costs every fusion about 7 points.}
\end{frame}

\begin{frame}{E2: class names masked}
  \takeaway{Masking halves the text model; then F2 is the best fusion, 2.5 points above F0.}

  \makebox[\textwidth]{%
  \slidetable
  \begin{tabular}{@{}lrrr@{}}
    \toprule
    \headrow
    macro-F1 & titles as is & names masked & gain over UB, masked \\
    \midrule
    UB text & $86.8$ & \F{$43.7$} & \\
    F0 late & $93.9$ & $78.9$ & $+35.2$ \\
    F1 concat & $92.8$ & $80.4$ & $+36.7$ \\
    F2 cross-modal & $93.5$ & \K{$81.3$} & $+37.6$ \\
    \bottomrule
  \end{tabular}}

  \showtakeaway
  \note{F2-mask vs F0-mask: +2.46, p = 2e-28. F2-mask vs F1-mask: +0.88, p = 2e-4. T0 masked: 47.8.
    On titles that never named the class, UB scores 39.2 vs 39.7 masked: the mask removes the name
    and nothing else. When no modality is decisive, token-level attention combines partial
    evidence better than averaging two probability vectors.}
\end{frame}

% ===========================================================================
\section{Analysis}
% ===========================================================================

\begin{frame}{Support and conflict}
  \takeaway{The gain is in the conflicts: F2 fixes more misleading titles than F1, F0 still the most.}

  \fig{support_conflict}

  \showtakeaway
  \note{Groups from UA and UB of the same seed, 3 seeds. Conflict-I = only the image model right
    (9 \% of items): F0 90.6, F2 75.3, F1 69.0. Conflict-T: all fusions about 97 to 98. Only the
    trained fusions solve some items where both unimodal models are wrong (F2 24 \%). Right panel,
    masked titles: Conflict-I grows to 40 \%, and F2-mask leads on Conflict-T (85.7 vs 75.9).}
\end{frame}

\begin{frame}{Inside F2}
  \takeaway{Uninformative titles get almost no attention, so F2 decides from the image.}

  \fig{qual_conflict_I}

  \showtakeaway
  \note{Conflict-I examples, seed 42. Grad-CAM of F2's image encoder for its prediction.
    [FUS] attention on image tokens: 0.91 on Conflict-I, 0.78 on Conflict-T. Compare between
    groups only: 49 image tokens vs about 15 title tokens, so uniform attention would give 0.75.}
\end{frame}

\begin{frame}{Why the unimodal models fail}
  \takeaway{Generic titles are easy to fix; a title that names another dish is not.}

  \makebox[\textwidth]{%
  \slidetable
  \begin{tabular}{@{}llrrrr@{}}
    \toprule
    \headrow
    failing model & main cause & items & F0 & F1 & F2 \\
    \midrule
    text  & title generic (``Hello'', site name) & $42$ & $100$ & $74$ & $93$ \\
    text  & title names another dish & $8$ & \F{$50$} & \F{$38$} & \F{$50$} \\
    image & looks like another dish & $34$ & $94$ & $94$ & $100$ \\
    image & image unrelated / wrong label / other & $16$ & $100$ & $100$ & $100$ \\
    \bottomrule
  \end{tabular}}

  \showtakeaway
  \note{100 random conflict items, seed 42, labelled by hand: Thanh 1-50, Phuoc 51-100.
    Example of the hard case: ``French Bread Pizza'' on a pizza photo, read as garlic bread.
    28 \% of the reviewed titles also occur in train: ``- quoteko.com'' is breakfast burrito in
    77 of 80 train items. Small sample: Wilson 95 \% CI of the 84 \% generic share is 71 to 92 \%.}
\end{frame}

\begin{frame}{Cost}
  \takeaway{F2's improvement is cheap: $+7\%$ compute over F1, same training time.}

  \makebox[\textwidth]{%
  \slidetable
  \begin{tabular}{@{}lrrrr@{}}
    \toprule
    \headrow
    & params & GMACs / item & train / run & latency (batch 1) \\
    \midrule
    UA & $23.7$\,M & $4.09$ & $14.4$ min & $4.7$ ms \\
    UB & $66.4$\,M & $0.64$ & $2.0$ min & $2.7$ ms \\
    F1 & $91.9$\,M & $4.73$ & $18.4$ min & $7.7$ ms \\
    F2 & $98.0$\,M & $5.06$ & $18.6$ min & $8.5$ ms \\
    \bottomrule
  \end{tabular}}

  \showtakeaway
  \note{RTX A6000 48 GB, shared lab machine, bf16. F0 costs UA + UB at inference and needs no
    extra training. Whole assignment: M2 2.0 GPU-hours, M3 about 5.7 GPU-hours.}
\end{frame}

\begin{frame}{Conclusions}
  \takeaway{Use F0 for informative titles; use F2 when the text is weak or confidence must be calibrated.}

  \begin{itemize}\setlength\itemsep{1.0em}
    \item fusion: $+6$ to $7$ points over the best single modality
    \item title: strong but leaky; check masked and unseen titles
    \item F2: beats F1, image branch kept alive, robust to a missing title
    \item F2 vs F0: tie on plain titles, \K{F2 $+2.5$} with names masked
    \item open problem: titles that name another dish
  \end{itemize}

  \showtakeaway
  \note{Next steps: a food-pretrained image encoder (A2), image-text contrastive pretraining, and
    a split with titles deduplicated across train and test.}
\end{frame}

\begin{frame}[plain]{}
  \takeaway{Fusion beats each modality by 6 to 7 points; the cross-modal F2 wins when the title stops naming the dish.}

  \makebox[\textwidth]{%
  \renewcommand{\arraystretch}{1.8}
  \begin{tabular}{@{}ll@{}}
    Code & \texttt{github.com/pbt245/CO3133-A3} \\
    Report and slides & \texttt{docs/a3-report/main.pdf}, \texttt{docs/a3-slides/slides.pdf} \\
    Checkpoints & \texttt{kaggle.com/datasets/bathanhphan/co3133-a3-m3-checkpoints} \\
    Group G-roundhog & Phan Bá Thanh (2353084), Trần Công Hoàng Phước (2352966) \\
  \end{tabular}}

  \showtakeaway
\end{frame}

\end{document}
